\section{TYPE-006: Loose Type for Constraint Options}
\label{sec:type-006}
\subsection*{Metadata}
\begin{tabular}{ll}
\textbf{Issue ID} & TYPE-006 \\
\textbf{Severity} & MEDIUM \\
\textbf{Category} & Type Safety \\
\textbf{Branch} & \branchref{fix/TYPE-006-type-constraint-options}
\textbf{Changes} & \branchdiff{fix/TYPE-006-type-constraint-options} \\
\end{tabular}

\subsection*{Executive Summary}
At \srcfile{src/sdk.ts:284}, the \code{constraints} option is cast as \code{any} before accessing its properties. This defeats TypeScript's checking for constraint field names and types. A misspelled field like \code{temprature} instead of \code{temperature} compiles without error but is silently ignored at runtime. The \code{as any} cast exists to support snake\_case variants (\code{max\_tokens}, \code{top\_p}, \code{top\_k}) that are not in the type definition.

\subsection*{Root Cause Analysis}
The \code{as any} cast bypasses type checking:

\begin{lstlisting}[language=TypeScript]
// BEFORE — as any cast to support snake_case
const c = constraints as any;
if (c.temperature !== undefined) modelOptions.temperature = c.temperature;
if (c.maxTokens !== undefined) modelOptions.maxTokens = c.maxTokens;
if (c.max_tokens !== undefined) modelOptions.maxTokens = c.max_tokens;
\end{lstlisting}

\subsection*{Fix Implementation}
Snake\_case variants are added to the type definition and the cast is removed:

\begin{lstlisting}[language=TypeScript]
// src/sdk-types.ts — add snake_case variants
constraints?: {
    temperature?: number;
    maxTokens?: number;
    max_tokens?: number;
    topP?: number;
    top_p?: number;
    topK?: number;
    top_k?: number;
};

// src/sdk.ts — remove as any cast
const c = constraints;  // No cast needed
if (c.temperature !== undefined) modelOptions.temperature = c.temperature;
if (c.maxTokens !== undefined) modelOptions.maxTokens = c.maxTokens;
if (c.max_tokens !== undefined) modelOptions.maxTokens = c.max_tokens;
if (c.topP !== undefined) modelOptions.topP = c.topP;
if (c.top_p !== undefined) modelOptions.topP = c.top_p;
if (c.topK !== undefined) modelOptions.topK = c.topK;
if (c.top_k !== undefined) modelOptions.topK = c.top_k;
\end{lstlisting}

\subsection*{Affected Files}
\begin{itemize}
    \item \srcfile{src/sdk-types.ts} — ConstraintOptions type
    \item \srcfile{src/sdk.ts} — Constraint mapping logic
    \item \srcfile{src/\_\_tests\_\_/type-006.test.ts}
\end{itemize}

\subsection*{Verification}
All constraint properties are now properly typed without \code{as any}.
